% Generated mechanically from the frozen stacks-introduction.tex target.
% Do not edit this generated file; edit the bound locale source instead.

% OK, start here.
%

\setcounter{chapter}{104}
\chapter{代数叠导论}


\phantomsection
\label{stacks-introduction-section-phantom}





\section{为什么阅读本章？}
\label{stacks-introduction-section-introduction}

\noindent
我们非正式地介绍代数叠。目标是迅速引入一套简单语言，
使你能够用它思考自己所关心的模问题的局部与整体性质。
做到这一点以后，在假定一般理论已经存在的前提下，你应当能够
对模问题提出适定的问题并开始求解。如果最终得到一个有趣的结果，
便可以回到 Stacks Project 其他部分的一般理论中，按需补齐细节。

\medskip\noindent
这里采用的观点接近 \cite{KatzMazur} 和 \cite{mumford_picard}
所采用的观点。






\section{预备知识}
\label{stacks-introduction-section-preliminary}

\noindent
设 $S$ 为概形。$S$ 上的一个{\it 椭圆曲线}是一个三元组
$(E, f, 0)$，其中 $E$ 是概形，且 $f : E \to S$ 和 $0 : S \to E$
是概形态射，并满足
\begin{enumerate}
\item $f : E \to S$ 是固有、光滑且相对维数为 $1$ 的态射，
\item 对每个 $s \in S$，纤维 $E_s$ 是亏格为 $1$ 的连通曲线，
即 $H^0(E_s, \mathcal{O})$ 和 $H^1(E_s, \mathcal{O})$
都是 $1$-维 $\kappa(s)$-向量空间，并且
\item $0$ 是 $f$ 的一个截面。
\end{enumerate}
给定椭圆曲线 $(E, f, 0)/S$ 和 $(E', f', 0')/S'$，
一个{\it 位于 $a : S \to S'$ 之上的椭圆曲线态射}
是态射 $\alpha : E \to E'$，使得图表
$$
\xymatrix{
E \ar[rr]_\alpha \ar[d]^f & & E' \ar[d]_{f'}  \\
S \ar@/^5ex/[u]^0 \ar[rr]^a & & S' \ar@/_5ex/[u]_{0'}
}
$$
交换且其中的方形为笛卡儿方形；换言之，态射 $\alpha$
诱导一个同构 $E \to S \times_{S'} E'$。
我们将定义椭圆曲线叠 $\mathcal{M}_{1, 1}$。
在 Stacks Project 的其余部分中，我们完整展开 Deligne 与 Mumford
在论文 \cite{DM} 中引入的方法：把 $\mathcal{M}_{1, 1}$ 表示为
一个带有函子的范畴
$$
p : \mathcal{M}_{1, 1} \longrightarrow \Sch, \quad
(E, f, 0)/S \longmapsto S
$$
% P12-ERR-0053: see ../control/ERRATA_CANDIDATES.jsonl
这意味着要使用概形范畴上的纤维化范畴、拓扑、群胚纤维化叠、
覆盖等等。在本章中，我们暂且把这些全部搁置，以稍微不同的方式
思考问题——这种方式或许更接近该理论创立者最初的思路。


\section{椭圆曲线的模叠}
\label{stacks-introduction-section-moduli-elliptic-curves}

\noindent
我们将做如下事情：
\begin{enumerate}
\item 从你所偏好的概形范畴 $\Sch$ 出发。
\item 加入一个新符号 $\mathcal{M}_{1, 1}$。
\item 一个态射 $S \to \mathcal{M}_{1, 1}$ {\bf 就是} $S$ 上的一条椭圆曲线
$(E, f, 0)$。
% P12-ERR-0054: see ../control/ERRATA_CANDIDATES.jsonl
\item 图表
$$
\xymatrix{
S \ar[rr]_a \ar[rd]_{(E, f, 0)} & & S' \ar[ld]^{(E', F', 0')} \\
& \mathcal{M}_{1, 1}
}
$$
{\bf 交换}，当且仅当存在一个位于 $a : S \to S'$ 之上的椭圆曲线态射
$\alpha : E \to E'$。我们称 $\alpha$ {\it 见证}了该图表的交换性。
\item 注意，交换图表按如下方式粘合：
$$
\xymatrix{
S \ar[rrr]_a \ar[rrrd]_{(E, f, 0)} & & &
S' \ar[d]_{(E', F', 0')} \ar[rrr]_{a'} & & &
S'' \ar[llld]^{(E'', F'', 0'')}
\\
& & & \mathcal{M}_{1, 1}
}
$$
因为若 $\alpha$ 和 $\alpha'$ 分别见证左、右三角形的交换性，
则 $\alpha' \circ \alpha$ 见证外侧三角形的交换性。
\item 复合
$$
S \xrightarrow{a} S' \xrightarrow{(E', f', 0')} \mathcal{M}_{1, 1}
$$
由 $(E' \times_{S'} S, f' \times_{S'} S, 0' \times_{S'} S)$ 给出。
\end{enumerate}
完成这一过程后，我们恰好用一个对象扩充了概形范畴 $\Sch$……

\medskip\noindent
不过，我们还没有定义从 $\mathcal{M}_{1, 1}$ 到概形 $T$ 的态射是什么。
答案是：采用使复合有意义的最弱概念。因此，态射
$F : \mathcal{M}_{1, 1} \to T$ 是一条规则，它为每条椭圆曲线
$(E, f, 0)/S$ 指定一个态射 $F(E, f, 0) : S \to T$，并且对任意交换图表
$$
\xymatrix{
S \ar[rr]_a \ar[rd]_{(E, f, 0)} & & S' \ar[ld]^{(E', F', 0')} \\
& \mathcal{M}_{1, 1}
}
$$
图表
$$
\xymatrix{
S \ar[rr]_a \ar[rd]_{F(E, f, 0)} & & S' \ar[ld]^{F(E', F', 0')} \\
& T
}
$$
也交换。你或许听说过的 $j$-不变量
$$
j : \mathcal{M}_{1, 1} \longrightarrow \mathbf{A}^1_{\mathbf{Z}}
$$
就是一个例子。这样看来，我们已经完成了……

\medskip\noindent
然而还没有！我们仍需定义态射
$\mathcal{M}_{1, 1} \to \mathcal{M}_{1, 1}$ 的概念。
做法与之前完全相同：态射
$F : \mathcal{M}_{1, 1} \to \mathcal{M}_{1, 1}$
是一条规则，它为每条椭圆曲线 $(E, f, 0)/S$ 指定另一条椭圆曲线
$F(E, f, 0)$，并保持上述图表的交换性。不过，由于我不知道
这种函子的非平凡例子，眼下就把从 $\mathcal{M}_{1, 1}$ 到自身的
态射集合定义为只含恒等态射。

\medskip\noindent
希望你已经看出怎样向这个扩充后的范畴加入其他对象。
直观上似乎很清楚：对任意“性态良好”的模问题，都可以执行上述构造，
向我们的范畴加入一个对象。事实上，现代代数几何的很大一部分
% P12-ERR-0055: see ../control/ERRATA_CANDIDATES.jsonl
正发生在这样一个世界中：$\Sch$ 被可数多个（明确构造的）模叠所扩充。

\medskip\noindent
你也许会反对说，由此得到的并不是范畴，因为我们必须给出交换性的见证，
所以对于图表何时交换、哪些图表的组合仍然交换，存在某种“含混性”。
然而事实证明，这种为交换性配备见证的想法正是处理 $2$-范畴的
一种有效方式！因此我们继续采用它。






\section{纤维积}
\label{stacks-introduction-section-fibre-products}

\noindent
这里要问的是，下面的纤维积应当是什么：
$$
\xymatrix{
& ? \ar@{..>}[rd] \ar@{..>}[ld] \\
S \ar[rd]_{(E, f, 0)} & & S' \ar[ld]^{(E', f', 0')} \\
& \mathcal{M}_{1, 1}
}
$$
答案是：从概形 $T$ 到 $? $ 的一个态射应当是三元组
$(a, a', \alpha)$，其中 $a : T \to S$、$a' : T \to S'$
是概形态射，而
$\alpha : E \times_{S, a} T \to E' \times_{S', a'} T$
是 $T$ 上椭圆曲线的一个同构。根据前面对复合与交换图表的定义，
这确实有意义。

\begin{lemma}[关键事实]
\label{stacks-introduction-lemma-key-fact}
函子 $\Sch^{opp} \to \textit{Sets}$，
$T \mapsto \{(a, a', \alpha)\text{ 如上所述}\}$
可由一个概形 $S \times_{\mathcal{M}_{1, 1}} S'$ 表示。
\end{lemma}

\begin{proof}
证明思路。把这个函子与
$$
\mathit{Isom}_{S \times S'}(E \times S', S \times E')
$$
联系起来，并使用 Grothendieck 的 Hilbert 概形理论。
\end{proof}

\begin{remark}
\label{stacks-introduction-remark-diagonal}
我们有公式
$S \times_{\mathcal{M}_{1, 1}} S' =
(S \times S')
\times_{\mathcal{M}_{1, 1} \times \mathcal{M}_{1, 1}}
\mathcal{M}_{1, 1}$。
因此，关键事实是 $\mathcal{M}_{1, 1}$ 的对角态射
$\Delta_{\mathcal{M}_{1, 1}}$ 的一个性质。
\end{remark}

\noindent
无论如何，关键事实使我们能够作如下定义。

\begin{definition}
\label{stacks-introduction-definition-smooth}
若对每个态射 $S' \to \mathcal{M}_{1, 1}$，投影态射
$$
S \times_{\mathcal{M}_{1, 1}} S' \longrightarrow S'
$$
都是光滑的，则称态射 $S \to \mathcal{M}_{1, 1}$ {\it 光滑}。
\end{definition}

\noindent
这与概形之间光滑态射的概念相容，因为光滑态射的基变换仍然光滑。
此外，很清楚怎样把这个定义推广到以 $\mathcal{M}_{1, 1}$
（或你所偏好的模叠）为靶的态射的其他性质。特别地，下面还会将它
用于{\it 满射}。





\section{定义}
\label{stacks-introduction-section-definition}

\noindent
我们把它表述为定义而不是结论，因为希望读者尝试其他情形
（不只是叠 $\mathcal{M}_{1, 1}$，也不只是全体概形的范畴 $\Sch$）。

\begin{definition}
\label{stacks-introduction-definition-algebraic-stack}
称 $\mathcal{M}_{1, 1}$ 为一个{\it 代数叠}，当且仅当
\begin{enumerate}
\item 对 $\Sch$ 上的 \'etale 拓扑，对象满足下降；
\item 关键事实成立；
% P12-ERR-0056: see ../control/ERRATA_CANDIDATES.jsonl
\item 存在满且光滑的态射 $S \to \mathcal{M}_{1, 1}$。
\end{enumerate}
\end{definition}

\noindent
第一个条件是一种“层性质”。这里把它详细写出，因为有一个技术点
需要说明。设给定概形 $S$、一个 \'etale 覆盖 $\{S_i \to S\}$，
以及态射 $e_i : S_i \to \mathcal{M}_{1, 1}$，使图表
$$
\xymatrix{
S_i \times_S S_j \ar[rd]_{e_i \circ \text{pr}_1} \ar[rr]_{\text{id}} & &
S_i \times_S S_j \ar[ld]^{e_j \circ \text{pr}_2} \\
& \mathcal{M}_{1, 1}
}
$$
交换。在这种情形下，层条件{\it 并不}保证存在态射
$e : S \to \mathcal{M}_{1, 1}$。事实上，我们还需为上述图表选取见证
$\alpha_{ij}$，并要求在 $S_i \times_S S_j \times_S S_k$ 上，作为见证有
$$
\text{pr}_{02}^*\alpha_{ik} =
\text{pr}_{12}^*\alpha_{jk} \circ \text{pr}_{01}^*\alpha_{ij}
$$
我想这是什么意思应当很清楚……若仍不清楚，恐怕就得阅读一些
关于群胚纤维化范畴等内容。无论如何，上式通常称为{\it 余圈条件}。
“层性质”的更精确表述是：给定 $\{S_i \to S\}$、
$e_i : S_i \to \mathcal{M}_{1, 1}$ 以及满足余圈条件的见证
$\alpha_{ij}$，存在一个在唯一同构意义下唯一的
$e : S \to \mathcal{M}_{1, 1}$，使得 $e_i \cong e|_{S_i}$，
并可恢复各 $\alpha_{ij}$。

\medskip\noindent
可以看到，即使精确表述这个命题也要花些工夫。
这种“层性质”的证明依赖于代数几何的一项基本技术，即下降理论。
我的建议是先直接接受“层性质”成立，再观察它在实际中蕴含什么。
事实上，要把“层性质”提炼成一个可以写在餐巾纸上的可操作命题，
需要相当的思维灵活性。也许最简单且已经颇有意思的变体如下：
设 $L/K$ 是域的有限伽罗瓦扩张，伽罗瓦群为
$G = \text{Gal}(L/K)$。令 $T = \Spec(L)$、$S = \Spec(K)$。
则 $\{T \to S\}$ 是一个 \'etale 覆盖。
令 $(E, f, 0)$ 为 $L$ 上的一条椭圆曲线。（是的，这只意味着
$E \subset \mathbf{P}^2_L$ 由一个 Weierstrass 方程给出，
而 $0$ 是通常的无穷远点。）记
$E_\sigma = E \times_{T, \Spec(\sigma)} T$ 为该基变换。
（是的，这对应于把 $\sigma$ 作用到 Weierstrass 方程的系数上，
还是应当用 $\sigma^{-1}$？）进一步设对每个 $\sigma \in G$，
给定 $T$ 上的同构
$$
\alpha_\sigma : E \longrightarrow E_\sigma
$$
此时，上述余圈条件意味着
$$
(\alpha_\tau)^\sigma \circ \alpha_\sigma = \alpha_{\tau\sigma}
$$
对所有 $\sigma, \tau \in G$ 成立。如果你学过群上同调，
这应当十分熟悉。无论如何，$\mathcal{M}_{1, 1}$ 上的“粘合”条件说，
若这一方程组有解，则存在 $S$ 上的一条椭圆曲线 $E'$，使得
$E \cong E' \times_S T$（实际上还说得更多，因为它也告诉你
怎样恢复各 $\alpha_\sigma$）。

\medskip\noindent
挑战：你能否完全只用 Weierstrass 方程所定义的椭圆曲线来证明这一点？




\section{一个光滑覆盖}
\label{stacks-introduction-section-smooth}

\noindent
最后要做的是找出 $\mathcal{M}_{1, 1}$ 的一个光滑覆盖。
事实上，在某种意义下，光滑覆盖的存在{\it 蕴含}
\footnote{这里稍有取巧，因为检验光滑性时必须证明与关键事实
非常接近的结论——毕竟光滑性是用纤维积定义的。好处是，
只需在纤维积一侧是你试图证明给出光滑覆盖的那个态射时，
证明这些纤维积存在。}关键事实！
对椭圆曲线，我们用 Weierstrass 方程来构造一个。

\medskip\noindent
令
$$
W = \Spec(\mathbf{Z}[a_1, a_2, a_3, a_4, a_6, 1/\Delta])
$$
其中 $\Delta \in \mathbf{Z}[a_1, a_2, a_3, a_4, a_6]$
是某个多项式（见下文）。令
$$
\mathbf{P}_W^2 \supset E_W :
zy^2 + a_1 xyz + a_3 yz^2 = x^3 + a_2x^2z + a_4xz^2 + a_6z^3.
$$
以 $f_W : E_W \to W$ 记投影。最后，以 $0_W : W \to E_W$
记由 $(0 : 1 : 0)$ 给出的 $f_W$ 的截面。事实证明，在赋权
$\deg(a_i) = i$ 下，存在一个关于 $a_i$ 的加权次数为 $12$ 的齐次多项式 $\Delta$，
使 $E_W \to W$ 光滑。可以通过计算 Weierstrass 方程的偏导数
显式求出它——当然也可以直接查表。还可以让 pari/gp 替你计算。
它就是
\begin{align*}
\Delta & = -a_6a_1^6 + a_4a_3a_1^5 + ((-a_3^2 - 12a_6)a_2 + a_4^2)a_1^4 + \\
& (8a_4a_3a_2 + (a_3^3 + 36a_6a_3))a_1^3 + \\
& ((-8a_3^2 - 48a_6)a_2^2 + 8a_4^2a_2 + (-30a_4a_3^2 + 72a_6a_4))a_1^2 + \\
& (16a_4a_3a_2^2 + (36a_3^3 + 144a_6a_3)a_2 - 96a_4^2a_3)a_1 + \\
& (-16a_3^2 - 64a_6)a_2^3 + 16a_4^2a_2^2 + (72a_4a_3^2 + 288a_6a_4)a_2 + \\
& -27a_3^4 - 216a_6a_3^2  -64a_4^3 - 432a_6^2
\end{align*}
你或许认得最后两项：在
$y^2 = x^3 + Ax + B$ 的情形，判别式为
$-64A^3 - 432B^2 = -16(4A^3 + 27B^2)$。

\begin{lemma}
\label{stacks-introduction-lemma-Weierstrass-smooth-cover}
态射 $W \xrightarrow{(E_W, f_W, 0_W)} \mathcal{M}_{1, 1}$
光滑且满。
\end{lemma}

\begin{proof}
满性来自每条域上的椭圆曲线都有 Weierstrass 方程这一事实。
% P12-ERR-0057: see ../control/ERRATA_CANDIDATES.jsonl
下面粗略说明一种证明光滑性的方法。考虑子群概形
$$
H =
\left\{
\left(
\begin{matrix}
u^2 & s & 0 \\
0 & u^3 & 0 \\
r & t & 1
\end{matrix}
\right)
\middle|
\begin{matrix}
u\text{ 为可逆元} \\
s, r, t\text{ 任意}
\end{matrix}
\right\}
\subset
\text{GL}_{3, \mathbf{Z}}
$$
$H$ 在 Weierstrass 概形 $W$ 上有作用 $H \times W \to W$。
为求出这一作用的方程，只需写出由 $H$ 中矩阵给出的坐标变换
如何作用于一般 Weierstrass 方程。随后可证明下列命题成立：
\begin{enumerate}
\item 任意椭圆曲线 $(E, f, 0)/S$ 在 $S$ 上 Zariski 局部地
都有一个 Weierstrass 方程；
\item $(E, f, 0)$ 的任意两个 Weierstrass 方程在 Zariski 局部
相差 $H$ 的一个元素。
\end{enumerate}
考察纤维积
$S \times_{\mathcal{M}_{1, 1}} W =
\mathit{Isom}_{S \times W}(E \times W, S \times E_W)$，
由此得出：这意味着态射 $W \to \mathcal{M}_{1, 1}$
是一个 $H$-挠子（即主齐性空间）。由于 $H \to \Spec(\mathbf{Z})$ 光滑，
而光滑群概形上的挠子也是光滑的，结论即得。
\end{proof}

\begin{remark}
\label{stacks-introduction-remark-quotient-stack}
上面概述的论证实际上表明
$\mathcal{M}_{1, 1} = [W/H]$ 是一个全局商叠。
证明一个模叠为代数叠的论证，约有 50\% 的时候还会表明
它是一个全局商叠。
\end{remark}



\section{代数叠的性质}
\label{stacks-introduction-section-properties}

\noindent
好，现在我们知道 $\mathcal{M}_{1, 1}$ 是代数叠了。
这能用来做什么？这里真正有帮助的，与其说是它为代数叠这一事实，
不如说是这样一种观点：$\mathcal{M}_{1, 1}$ 的性质应编码在态射
$S \to \mathcal{M}_{1, 1}$ 的性质中，也就是编码在椭圆曲线族中。
下面列出若干例子。

\medskip\noindent
{\bf 局部性质：}
$$
\mathcal{M}_{1, 1} \to \Spec(\mathbf{Z})\text{ 是光滑的}
\Leftrightarrow
W \to \Spec(\mathbf{Z})\text{ 是光滑的}
$$
{\bf 思路}。代数叠的局部性质编码在其光滑覆盖的局部性质中。

\medskip\noindent
{\bf 整体性质：}
$$
\begin{matrix}
\mathcal{M}_{1, 1}\text{ 是拟紧的} \Leftarrow W\text{ 是拟紧的}
\\
\mathcal{M}_{1, 1}\text{ 是不可约的} \Leftarrow W\text{ 是不可约的}
\end{matrix}
$$
{\bf 思路}。代数叠的某些整体性质可以从其一个{\it 合适的}
% P12-ERR-0058: see ../control/ERRATA_CANDIDATES.jsonl
\footnote{我猜想可能存在没有不可约光滑覆盖的不可约代数叠——
但如果确实存在，它一定相当糟糕！}光滑覆盖的相应性质读出。

\medskip\noindent
{\bf 拟凝聚层：}
$$
\QCoh(\mathcal{O}_{\mathcal{M}_{1, 1}}) =
H\text{-等变拟凝聚模，定义于 }W
$$
{\bf 思路}。一方面，$\mathcal{M}_{1, 1}$ 上的一个拟凝聚模
应当对应于：对每个态射 $e : S \to \mathcal{M}_{1, 1}$，
在 $S$ 上给定一个拟凝聚层 $\mathcal{F}_{S, e}$。
特别地，可取态射 $(E_W, f_W, 0_W) :  W \to \mathcal{M}_{1, 1}$。
由于这个态射是 $H$-等变的，所得拟凝聚模 $\mathcal{F}_W$
也是 $H$-等变的。反过来，给定一个 $H$-等变模，从
$S \times_{e, \mathcal{M}_{1, 1}} W$ 是 $H$-挠子这一观察出发，
% P12-ERR-0059: see ../control/ERRATA_CANDIDATES.jsonl
可以通过下降理论恢复各层 $\mathcal{F}_{S, e}$。

\medskip\noindent
{\bf Picard 群：}
$$
\Pic(\mathcal{M}_{1, 1}) = \Pic_H(W) = \mathbf{Z}/12\mathbf{Z}
$$
{\bf 思路}。第一个等式已在上面看到。注意 $\Pic(W) = 0$，
因为环 $\mathbf{Z}[a_1, a_2, a_3, a_4, a_6, 1/\Delta]$
的类群平凡。存在正合序列
$$
\mathbf{Z}\Delta \to
\Pic_H(\mathbf{A}^5_{\mathbf{Z}}) \to
\Pic_H(W) \to 0
$$
中间的群等于 $\Hom(H, \mathbf{G}_m) = \mathbf{Z}$。
% P12-ERR-0060: see ../control/ERRATA_CANDIDATES.jsonl
$\Delta$ 的像为 $12$，因为 $\Delta$ 的次数是 $12$。
这个论证大致正确，见 \cite{PicM11}。

\medskip\noindent
{\bf \'Etale 上同调：}设 $\Lambda$ 为环。
% P12-ERR-0061: see ../control/ERRATA_CANDIDATES.jsonl
存在一个第一象限谱序列，收敛到
$H^{p + q}_\etale(\mathcal{M}_{1, 1}, \Lambda)$，其 $E_2$-页为
$$
E_2^{p, q} = H_\etale^q(W \times H \times \ldots \times H, \Lambda)
\quad(p\text{ 个因子，每个均为 }H)
$$
{\bf 思路}。注意
$$
W \times_{\mathcal{M}_{1, 1}} W \times_{\mathcal{M}_{1, 1}}
\ldots \times_{\mathcal{M}_{1, 1}} W
= W \times H \times \ldots \times H
$$
因为 $W \to \mathcal{M}_{1, 1}$ 是一个 $H$-挠子。
这个谱序列是光滑覆盖 $\{W \to \mathcal{M}_{1, 1}\}$ 的
{\v C}ech-上同调谱序列。例如，由于 $W$ 连通，我们看到
$H^0_\etale(\mathcal{M}_{1, 1}, \Lambda) = \Lambda$；
又因为 $H^1_\etale(W, \Lambda) = 0$，所以
$H^1_\etale(\mathcal{M}_{1, 1}, \Lambda) = 0$
（当然，这需要证明）。当然，光滑覆盖
$W \to \mathcal{M}_{1, 1}$ 对计算 \'etale 上同调而言未必“最优”。
